\section{Part E. Categories of curves and the arithmetic image}

The manuscript loci are cote 142, manuscript pp.~7--10, scans
15--21. We have opened the supplied images of scans 17 and 19.
The phrase ``programme \`a l'horizon'' on scan 17 concerns the
arithmetic action as well as the category. On scan 19 the proposed
isotopy relation is explicitly tentative. Neither passage licenses
identifying an arbitrary path of local homeomorphisms with a
two-sided orbit. We distinguish full \emph{\'etale} paths in a moduli
stack, paths tame at its compactification boundary, and its maximal
prime-to-$p$ quotient. These are three different objects.

\subsection{Two obstructions to the requested positive-characteristic construction}

\begin{claim}{E.1}{Proved; an obstruction with explicit hypotheses}
Let $k=\overline{\F}_p$ and $k'=\overline{k(t)}$. On
$U=\mathbb P^1_k\setminus\{0,1,\infty\}$ the connected
$\Z/p$-torsor
\[
                       z^p-z=tx
\]
over $U_{k'}$ does not descend to $U$, even after forgetting its
chosen identification of the deck group with $\Z/p$.
Consequently the base-extension map of full geometric fundamental
groups is not an isomorphism. A category having the marking
groupoids of the full moduli stacks as its isomorphism groupoids
cannot have base extension as an equivalence in characteristic $p$.
This remains an obstruction if only the noninvertible covering
arrows are restricted to prime-to-$p$ Galois closures.
\end{claim}
\begin{proof}
The Artin--Schreier sequence on the affine curve gives
\[
H^1(U_{k'},\Z/p)=
k'[x,x^{-1},(x-1)^{-1}]/\{h^p-h\}.
\]
Modulo an Artin--Schreier coboundary, a principal part has a unique
representative whose pole orders are prime to $p$, after successively
removing terms of order divisible by $p$. The field is perfect, so
the necessary $p$-th roots of coefficients exist uniquely. Uniqueness
follows because a nonconstant $h$ has at every pole a leading term
of $h^p-h$ of order divisible by $p$. Constants contribute nothing
over an algebraically closed field. In the displayed class the sole
pole is at infinity, of order one, with coefficient $t$. A class
defined over $k$ has reduced principal parts with coefficients in
$k$. Therefore this class cannot descend. It is nonzero, which makes
the torsor connected. Relabelling its deck group multiplies its
class by an element of $\F_p^*$ and does not remove $t$.

Now $\mathcal M_{0,4}$ with ordered marks is this same curve.
Its full fundamental group is therefore not invariant under this
base extension. For unordered marks, the ordered moduli scheme is
a finite \emph{\'etale} cover of the quotient stack, with constant
permutation group $S_4$. An isomorphism on the fundamental groups
of the unordered stack, compatible with this cover, would induce
an isomorphism on the corresponding open kernels, contrary to
the ordered calculation. Thus either convention has the obstruction.

In the factorization category of (Z.2), an endomorphism of a fixed
hyperbolic type which factors as a marking followed by a finite
\emph{\'etale} cover must have covering degree one. Indeed, for
a degree-$d$ cover of open curves in characteristic $p$, the
Riemann--Hurwitz formula, including the nonnegative wild different
at the boundary, gives
\[
        2g_X-2+n_X\ \geq\ d(2g_Y-2+n_Y).
\]
Equality of the two hyperbolic types forces $d=1$. The isomorphism
groupoid is consequently precisely the specified marking groupoid,
when markings are retained faithfully as in (Z.2). Its failure of
base invariance already contradicts equivalence of categories.

The qualification about faithful markings matters. If an unspecified
``category'' is allowed to identify arbitrary paths or arbitrary
covering arrows, the words ``arrows include'' alone do not determine
its Hom sets. The Artin--Schreier example is not a proof about every
such artificial quotient category. It proves the obstruction for
the exact full-path construction requested by the manuscript and
R26. The coefficient and pole-order calculations for $p=2,3,5$ are
also checked in \texttt{code/root\_checks.py}; the argument above
proves it for every $p$.
\end{proof}

\begin{claim}{E.2}{Refuted}
There is no type-preserving functor from the entire characteristic-zero
category $\mathcal T$ to a category whose arrows all factor as a
type-preserving marking followed by a finite \emph{\'etale} map
with prime-to-$p$ Galois closure, if the functor is required to
specialize every finite-cover arrow. This obstruction holds for
every prime $p$, including $2$ and $3$.
\end{claim}
\begin{proof}
Let $C/\C$ be a smooth proper genus-two curve. The presentation
\[
\pi_1(C^{\rm an})=
\langle a_1,b_1,a_2,b_2\mid[a_1,b_1][a_2,b_2]=1\rangle
\]
surjects onto $\Z/p$ by sending $a_1$ to one and the other
generators to zero. The connected unramified topological cover
algebraizes by SGA~1 XII, Theorem~5.1 and Corollary~5.2
\citep{E_SGA1}. Its source $D$ has genus $p+1$, because
$2g_D-2=p(2g_C-2)=2p$.

A type-preserving functor would give an arrow from type $(p+1,0)$
to type $(2,0)$. A marking does not change the type. Any subsequent
finite \emph{\'etale} cover between these proper types in
characteristic $p$ has degree
\[
                  \frac{2(p+1)-2}{2\cdot2-2}=p.
\]
The order of its geometric Galois closure is divisible by its degree,
so cannot be prime to $p$. No such arrow exists. This contradiction
does not depend on selecting a bad model for $D$: it occurs in the
types before any model is chosen.

Accordingly the source of a possible specialization functor must
at least be restricted to arrows whose geometric Galois closures
have order prime to $p$. Replacing all moduli markings by their
maximal prime-to-$p$ quotients is not automatically the right further
repair: a prime-to-$p$ covering group can have an automorphism group
with $p$-torsion, and transport of such covers can detect that
torsion. Tameness at the boundary is a different restriction.
\end{proof}

\begin{claim}{E.3}{Known, with derivation; the fixed-curve replacement}
Let $(\mathcal C,\mathcal D)$ be a smooth proper curve with finite
\emph{\'etale} marked divisor over a strictly henselian complete
discrete valuation ring of mixed characteristic $(0,p)$, with
algebraically closed residue field. Specialization induces an
isomorphism on the maximal prime-to-$p$ geometric fundamental
groups of the two open fibres. Hence connected covers with
prime-to-$p$ Galois closure correspond, preserving degree, deck
group, and the type of source. The analogous category of covers
of a fixed smooth marked curve over an algebraically closed
characteristic-$p$ field is invariant under extension of that field.
\end{claim}
\begin{proof}
The precise sources are SGA~1 XIII, Corollary~2.12 and
Proposition~4.3, its Examples~4.4, and the base-change statement
4.5.3 \citep{E_SGA1}. Their hypotheses are a proper smooth
relative curve, a relatively normal-crossings boundary finite
\emph{\'etale} over the base, and orders invertible at the
specialization. A section exists by lifting a smooth residue point
over the strictly henselian base. These hypotheses are satisfied here. In this setting
specialization is surjective on the tame group and is an isomorphism
on its maximal prime-to-residue-characteristic quotient. A connected
cover with prime-to-$p$ Galois closure is a transitive set for a
finite prime-to-$p$ quotient of that group. An isomorphism of
the groups identifies these finite sets, their equivariant maps,
their sizes, and their automorphism groups. Boundary inertia
subgroups correspond by the same Kummer-parameter comparison;
thus their ramification indices and the number of points above
each mark are unchanged. Tame Riemann--Hurwitz then gives the same
genus. For an extension of algebraically closed fields use the
relative-group base-change statement in XIII, 4.5.3, applied over
the algebraically closed base field with a chosen point; the same
finite-set argument gives the asserted equivalence.

This is a theorem about covering categories over a fixed curve.
It does not prove the missing assertion about paths and composition
over all moduli stacks. Bertin--Romagny, Theorem~6.9, constructs
the smooth proper stable Hurwitz stack over $\Z[1/|G|]$;
Proposition~6.14 gives its proper quasi-finite flat discriminant
map, \emph{\'etale} over the smooth target locus; and Theorem~6.20
identifies the non-Galois-cover stack by the quotient recording
its Galois closure \citep{E_BR2007}. These are the exact additional
inputs available for that construction. They do not by themselves
say that the \emph{full} path groupoid of the nonproper moduli
stack is invariant; E.1 proves that it is not.
\end{proof}

\subsection{The arithmetic-image route and a theorem that it omits}

\begin{claim}{E.4}{Known, with derivation}
The conjugacy class of the image of $G_\Q$ in the profinite
Grothendieck--Teichm\"uller group has a published combinatorial
characterization. It is the class of maximal closed subgroups
satisfying the AA-property of Hoshi--Mochizuki--Tsujimura,
Definition~5.12, with the COF, RGC and ArBC hypotheses required
there. Determining this conjugacy class is therefore distinct
from proving the equality $\widehat{\mathrm{GT}}=G_\Q$.
\end{claim}
\begin{proof}
Hoshi--Mochizuki--Tsujimura, Theorem~5.17(ii), states exactly this
characterization; its input is the abstract profinite group of
the second configuration space of the tripod
\citep{E_HMT2025}. Definition~5.12 imposes four conditions on
the reconstructed Kummer and function-field data; AA includes
the second-configuration compatibility condition (iv), which
cannot be dropped. Theorem~5.15 puts every such subgroup inside
a conjugate of $G_\Q$, and Remark~5.12.1 verifies the conditions
for $G_\Q$ itself.

Here is the maximality step. In a profinite group $G$, if a closed
subgroup $H$ satisfies $H\subseteq gHg^{-1}$, then $H=gHg^{-1}$.
In every finite quotient the two images have equal cardinalities,
so the inclusion is equality; intersecting the inverse images over
all finite quotients gives equality of the closed subgroups.
Thus every AA subgroup is contained in an AA subgroup conjugate
to $G_\Q$, and no such conjugate is properly contained in another.
This proves the stated maximal-subgroup characterization from
the cited theorem, keeping its hypotheses intact.

This use does not assert that every GT element is arithmetic, or
that testing only its cyclotomic component and the pentagon is
sufficient. Nor does it yet construct the compatible action on
every type and every puncture-filling arrow asked for in E2.
\end{proof}

\subsection{An actual isotopy congruence and its profinite envelope}

We give a precise category before asking for its algebraic comparison.
Objects are oriented finite-type topological surfaces $S$ with
$c(S)=2g(S)-2+n(S)>0$. An elementary arrow is an oriented finite
unramified covering or an oriented puncture filling. A puncture
filling extends to an orientation-preserving homeomorphism of the
compact surfaces and fills a specified subset of the finite end set.
Homeomorphisms are included as degree-one arrows. A word is a
finite composable string of elementary arrows. It determines a
local homeomorphism on the open surfaces and a finite branched
cover on their end compactifications.

We impose the following equivalence, called \emph{diagram isotopy}.
An elementary isotopy is a continuous one-parameter family of a
fixed finite diagram of elementary arrows, with fixed endpoint
surfaces and fixed types at each vertex. We also identify words
having the same actual composite map. We take the equivalence
relation generated by these identifications and by their occurrence
inside a longer composable word. In particular this is explicitly
a congruence, rather than a relation whose compatibility with
composition is presumed. Denote the resulting category by
$\mathcal C_{\rm diag}$.

\begin{claim}{E.5}{Proved}
Composition is well-defined on $\mathcal C_{\rm diag}$.
The category of finite-cover isotopy classes embeds faithfully in
it, and its automorphism group at $S$ is the oriented mapping
class group $M_S$. It contains puncture fillings and the usual
forgetting functors on groupoids of surfaces with specified
labelled punctures. Every Hom set has finitely many right
$M_S$-orbits. In particular, for every arrow $u:S\to T$, its
target lifting subgroup
\[
 H_u=\{b\in M_T:\text{ there is }a\in M_S
                          \text{ with }bu=ua\}
\]
has finite index in $M_T$.
\end{claim}
\begin{proof}
Concatenation of words is associative, and the stipulated congruence
makes it independent of representatives. The degree of the compact
map is a positive integer. It is constant in every elementary
isotopy and multiplicative under composition; all relations preserve
it. If a word of degree $d$ joins $S$ to $T$, then
\[
                  c(S)=d\,c(T)+\Delta,
                       \qquad \Delta\in\Z_{\geq0}.
\]
For a covering, $\Delta=0$ by Euler characteristic. For a filling
of $r>0$ punctures the degree is one and $\Delta=r$.
For composites the formula is
$\Delta(vu)=\Delta(u)+\deg(u)\Delta(v)$, so every genuine
filling contributes positively. Equivalently $\Delta$ is the sum
of the local degrees at extra omitted source points whose images
are interior target points; this follows by Riemann--Hurwitz on
the compactifications.

If a word represents a finite covering, its endpoint type and
degree force $\Delta=0$. Every word equivalent to it also has
$\Delta=0$, so all its elementary factors are coverings.
The elementary isotopies then are ordinary covering isotopies,
and replacing a factorization by its composite makes no additional
identifications. Thus the embedding is faithful. An invertible
arrow has an inverse of positive degree, so both degrees are one.
The two complexity formulas force both defects to vanish; all
factors are homeomorphisms. Their isotopy classes form $M_S$.

Each noninvertible elementary arrow strictly decreases the positive
integer $c$. After absorbing adjacent homeomorphisms, the length
of a word from $S$ to $T$ is at most $c(S)-c(T)$, and its degrees
are bounded by $c(S)/c(T)$. Only finitely many hyperbolic types
have complexity at most $c(S)$. For a fixed target and degree
there are only finitely many connected covering types: a finitely
generated surface group has only finitely many homomorphisms to
the finite symmetric group of that degree. Each such covering
type gives one right source-mapping-class orbit.

For a filling, moving the finitely many extra points by an ambient
isotopy of the target identifies any two configurations of the same
cardinality. The isotopy may fix the retained punctures. The
remaining difference is a source homeomorphism. Thus fillings
also have finitely many right source orbits; for unlabelled extra
points there is one. This assertion includes the braid motions
of the extra points: their lifts act in the source stabilizer,
rather than producing an omitted family of arrows.

The property of having finitely many right orbits is preserved
under composition. Choose finite sets of representatives $u_i$
and $v_j$ for two adjacent Hom sets. Each composite has the
form $v_j b u_i a$. Left multiplication by $b$ permutes the
finite set of right orbits of the first Hom set. Therefore
finitely many further representatives suffice for all composites.
There are only finitely many possible elementary type-and-degree
patterns, proving the assertion for all words and hence for their
quotient. Finally $H_u$ is the stabilizer of the right orbit of
$u$ under the left action of $M_T$ on this finite set. Its index
is finite.

For labelled punctures, a homeomorphism preserving all labels
extends over any specified subset of punctures. Isotopies extend
in the same way. This defines the forgetting functor on the
isomorphism groupoids; successive forgettings equal forgetting
the union. The remaining type must still be hyperbolic. These
are the topological functors that the moduli-stack maps in R26
are to realize.
\end{proof}

\begin{claim}{E.6}{Proved}
There is a Hausdorff profinite envelope of
$\mathcal C_{\rm diag}$ with the same objects and
\[
\Hom_{\widehat{\mathcal C}_{\rm diag}}(S,T)
 =\coprod_{uM_S\in\Hom_{\mathcal C_{\rm diag}}(S,T)/M_S}
          \overline{S_u}\backslash\widehat M_S,
 \qquad S_u=\{a\in M_S:ua=u\}.
\]
Here the closure is in $\widehat M_S$. Each Hom space is compact,
composition is continuous, the original category has dense image,
and $\Aut(S)=\widehat M_S$. This proves the topological
completion part of E3 for the specified congruence.
\end{claim}
\begin{proof}
We verify the construction rather than invoke F149-G6 without
checking composition. Complete a right orbit $S_u\backslash M_S$
as the inverse limit of finite sets $S_uN\backslash M_S$ for
finite-index normal $N\lhd M_S$. Compactness identifies that
limit with the displayed closed-coset space. A target element
$b\in H_u$ has a lift $a$ with $bu=ua$. Its class in
$N_{M_S}(S_u)/S_u$ is well-defined and depends homomorphically
on $b$. Left multiplication by it acts on every finite set
$S_uN\backslash M_S$, through a finite quotient of $H_u$.
Since $H_u$ has finite index, its full profinite completion is
its closure in $\widehat M_T$: the normal cores in $M_T$ of
its finite-index subgroups give the required cofinality.
Thus its action extends continuously. The finitely many cosets
of $H_u$ extend the action to all of $\widehat M_T$, permuting
the finitely many right orbits in the double orbit of $u$.

For an original arrow $v:T\to V$, define
$L_v(u\alpha)=(vu)\alpha$. It is well-defined and continuous
because $S_u\subset S_{vu}$. For completed arrows set
\[
              (v\beta)(u\alpha)=L_v(\beta\cdot u)\alpha.
\]
Replacing $\alpha$ by its left $\overline{S_u}$-coset makes
no change by equivariance. Replacing $\beta$ by its left
$\overline{S_v}$-coset makes no change because
$L_v(sh)=L_v(h)$ first for original $s,h$, then for their
closures by continuity and density. Changes of right-orbit
representatives are translations of the same formula. On each
product of right-orbit pieces the formula is the quotient of a
continuous map of compact profinite groups, so is jointly
continuous. Associativity follows because the two continuous
triple-composition maps agree on the dense set of original
triples and the target is Hausdorff.

The orbit of the identity is $\widehat M_S$ and is open.
If two completed arrows are inverse, approximate them by original
arrows within their original orbit pieces and with both composites
in the open identity orbits. Those original composites are
automorphisms. Each of the approximating arrows, and hence its
orbit representative, is an isomorphism. This proves that no
new invertible orbit appears. It proves the asserted automorphism
group and conservativity. Compactness follows from the finiteness
of the right-orbit set in E.5. Any functor to a Hausdorff
topological category with profinite automorphism groups extends
uniquely: extend the homomorphisms on $M_S$ to $\widehat M_S$,
then the continuous orbit maps kill $\overline{S_u}$.
Compatibility with composition again follows by density. This is
the Hausdorff universal property required in R25.
\end{proof}

\begin{claim}{E.7}{Proved; the moving-critical-value test}
Diagram isotopy identifies the maps $f_0$ and $f_t$ in R26,
for $0\leq t\leq1/10$, but they are not in one two-sided
orbit for $t>0$. Thus the congruence of E.5 does not silently
repeat the false two-sided-orbit assertion.
\end{claim}
\begin{proof}
Use the four-factor diagram
\[
X_t\xrightarrow{z^3-3z+t}Y_{1,t}
 \lhook\joinrel\longrightarrow Y_2
 \xrightarrow{w^2}Y_3\lhook\joinrel\longrightarrow Y,
\]
where
\begin{align*}
Y&=\mathbb P^1\setminus\{9,16,\infty\},
&Y_3&=Y\setminus\{0\},\\
Y_2&=\mathbb P^1\setminus\{0,3,-3,4,-4,\infty\},
&Y_{1,t}&=Y_2\setminus\{t-2,t+2\},
\end{align*}
and $X_t$ is the inverse image of $Y_{1,t}$ under the cubic.
The cubic critical values are $t\pm2$, always distinct from
the fixed omitted values on this interval. The five relevant
discriminants are $27(4-(t-c)^2)$ for
$c=0,\pm3,\pm4$, and never vanish. The source has twenty
distinct punctures throughout. Ambient isotopies identifying
these moving punctures give a family of the entire fixed diagram,
so E.5 identifies its endpoints. On compactifications the finite
critical values of the composite are
$0,(t-2)^2,(t+2)^2$. Their number is two at zero and three at
positive $t$, since their difference is $8t$. A two-sided
homeomorphism preserves that number. The factorization and
discriminants are verified exactly in \texttt{code/root\_checks.py}.
\end{proof}

% Prepared by the Part F agent; independently reviewed by root.
\subsection{The algebraic Hausdorff presentation and its comparison}

We specify the algebraic category, including its separation convention.
Its objects are smooth connected hyperbolic curves over $\C$.
Use the stacks $\mathcal M_{g,[n]}$ of unordered punctures;
labelled or coloured variants use the corresponding finite covers.
Generators are paths in their \'etale fundamental groupoids,
actual finite \'etale maps, and actual puncture fillings.
An actual isomorphism is identified with its canonical stack path.
Besides the path-groupoid laws, impose the following relations.
For every finite chain of specified covers and fillings, let
$\mathcal H_w$ be its diagram stack. A path in $\mathcal H_w$
from $w_0$ to $w_1$, inducing source path $\alpha$ and target
path $\beta$, imposes
\[
                       \beta w_0=w_1\alpha.
\]
Two chains of actual algebraic maps with the same composite
are also identified. All these relations may occur inside
longer words.

The \emph{Hausdorff \'etale presentation} means the universal
Hausdorff topological category with profinite automorphism groups
realizing these data, continuously on each \'etale path groupoid.
This is a universal property in terms of algebraic stacks and
their \'etale paths alone. Equivalently, relations are imposed
in a closed congruence. The following proof establishes its
existence, so existence is not an extra assumption in the definition.

\begin{claim}{E.8}{Proved}
The Hausdorff \'etale presentation exists. After choices of
surface markings, it is equivalent to
$\widehat{\mathcal C}_{\rm diag}$ of E.6. In particular,
\[
 \operatorname{Hom}(S,T)=
 \coprod_{uM_S\in\operatorname{Hom}_{\mathcal C_{\rm diag}}(S,T)/M_S}
       \overline{S_u}\backslash\widehat M_S,
 \qquad S_u=\{a\in M_S:ua=u\}.
\]
Composition is continuous, Hom spaces are compact, and
$\operatorname{Aut}(S)=\widehat M_S$. The completed pure-cover
category embeds faithfully, and the labelled puncture-forgetting
operations agree with those induced by the algebraic moduli maps.
This proves the comparison and completion required in E3 for
the isotopy relation of E.5.
\end{claim}
\begin{proof}
First the diagram stacks used in the definition exist and are
of finite type. A filling is described by a configuration of
extra disjoint points, with the retained and forgotten sets
distinguished. For a covering use the finitely many smooth
Hurwitz stacks of the given degree and boundary ramification
partitions. Non-Galois maps are included by
\citet[Theorem~6.20]{E_BR2007}; its opened proof identifies
the local deformation of a smooth cover with that of its
marked target. Iterated fibre products over the intermediate
moduli stacks give $\mathcal H_w$. Hyperbolicity makes the
automorphism groups finite. We retain these stack automorphisms:
a Hurwitz map to target moduli need not be representable when
deck automorphisms occur, as the footnote to
\citet[Proposition~6.14]{E_BR2007} explicitly records.

\emph{The Betti presentation.}
Temporarily use ordinary paths in the analytic stacks.
Contractibility of Teichm\"uller space and its quotient
description identify the fundamental groupoid of each moduli
stack with the groupoid of oriented surface markings up to
isotopy \citep[\S2.1]{F_Schneps2003}.
An analytic diagram path, after local topological trivializations
over an interval, gives a diagram isotopy with its endpoint
transports. Equal algebraic composites give equal underlying
maps. These observations define a functor from the Betti
presentation to $\mathcal C_{\rm diag}$.

Conversely, choose a complex structure on each surface object.
For a topological word ending in $T$, keep its chosen target
structure and propagate that structure backwards: a covering
pulls it back, and a filling restricts it to the punctured open.
All factors become holomorphic, extend over their compactifications,
and are algebraic. Join the resulting source structure to its
chosen structure by a path of marked complex structures.
Contractibility of Teichm\"uller space makes the path with
this prescribed identity marking unique up to the relevant
homotopy.

Here is the compatibility with the relations, including the
family issue. A topological diagram isotopy can be represented
at the level of its paths and endpoint transports by smooth
finite-configuration paths and finite-cover transports. For
fillings this is ordinary approximation of paths in the open
configuration manifold, relative to the endpoints, followed
by isotopy extension. Covering paths are uniquely transported
after their starting cover is fixed. Work backwards from the
target: these configuration paths and covering transports
give a path in the analytic diagram stack. This asserts a
continuous stack path; it does not assert that an arbitrary
nonsmooth original isotopy is a holomorphic family in its
real parameter. The induced source path has the prescribed
identity marking, so the naturality relation gives the
required equality of Betti arrows.

If two words have the same actual composite, their pulled-back
source complex structures are the same. Both factorizations
are simultaneously holomorphic and have equal algebraic
composites. The stipulated relation therefore identifies them.
For composition, transport the first diagram along the marked
complex-structure path used for the second; its naturality
square identifies this construction with that of the concatenated
word. These arguments prove independence of the choices and
give inverse functors, up to marking isomorphisms. The Betti
presentation is therefore equivalent to $\mathcal C_{\rm diag}$.
No two-sided-orbit assertion is used; E.7 is still a diagram
path relation.

\emph{All \'etale path relations.}
Let $\mathcal A$ be any Hausdorff topological category with
profinite automorphism groups, and let
$F:\mathcal C_{\rm diag}\to\mathcal A$ be a functor.
By the Betti comparison this is a realization of the Betti
presentation. Each induced homomorphism from a moduli
fundamental group to an automorphism group of $\mathcal A$
extends uniquely to the full profinite completion. Choosing
one ordinary path to each endpoint extends this assertion
to path groupoids; changing the choice changes the formula
by an ordinary loop and leaves the extension unchanged.

Noohi's Theorem~20.4 and its following corollary identify
these completions with the \'etale fundamental groups
\citep{E_Noohi2005}. The finite-type complex Deligne--Mumford
stacks in question have locally path connected, semilocally
simply connected analytic stacks. Thus ordinary path torsors
are dense in the \'etale path torsors, component by component.
For a path $\gamma$ in any one $\mathcal H_w$, the two sides
of its naturality equation are continuous functions of
$\gamma$ into a Hausdorff Hom space of $\mathcal A$.
They agree on ordinary paths, hence on their closure.
The equality therefore holds for every \'etale path, for
every diagram stack, and after every insertion into a
longer word. The equal-actual-composite relations do not
change. Conversely restricting a continuous \'etale
realization to ordinary paths gives the original functor.
The two constructions are inverse, also on compatible
natural transformations.

\emph{Existence and separation.}
By E.6, $\widehat{\mathcal C}_{\rm diag}$ represents functors
from $\mathcal C_{\rm diag}$ to every such $\mathcal A$.
The preceding paragraph identifies this universal property
with that of the algebraic Hausdorff \'etale presentation.
It proves both existence and the asserted comparison.

We make the stabilizer conclusion explicit. Every Hausdorff
realization kills $\overline{S_u}$, since its orbit map is
continuous and kills $S_u$. Conversely, if
$a\notin\overline{S_u}$, the category of E.6 is itself
one realization of the algebraic presentation, by the
preceding construction. Its distinct cosets
$\overline{S_u}$ and $\overline{S_u}a$ separate $u$
from $ua$. No universal closed relation can therefore
identify them. This proves the exact closure formula
without a subgroup-separability hypothesis or an
interchange of completion with an ordinary quotient.

The remaining structural statements follow from E.6.
For a pure covering, defect zero excludes all genuine
fillings at every intermediate relation by E.5. Its
discrete stabilizers are the covering-only ones, and so
are their closures. This proves the asserted completed
embedding. Labelled forgetting maps and their composites
are included in the presentation by the literal
algebraic forgetting maps and their equality relations.
\end{proof}

\begin{claim}{E.9}{Proved}
For the degree-two cover followed by one filling
\[
 f_b(z)=z^2+b,\qquad b\in Y:=\mathbb P^1\setminus\{0,1,\infty\},
\]
\[
 X_b=\mathbb P^1\setminus
 \{0,\infty,\sqrt{-b},-\sqrt{-b},\sqrt{1-b},-\sqrt{1-b}\},
\]
fix the three labelled target punctures. The diagram component
has stack $Y\times B\mu_2$. Its source braid subgroup is
the image of
\[
 \phi:\pi_1(Y)\times C_2\longrightarrow M_{0,[6]}.
\]
This is the right stabilizer of the arrow in the indicated
ramification component. Its completed stabilizer is exactly
$\overline{\operatorname{im}\phi}$, equal to the image
of $\pi_1^{\mathrm{et}}(Y\times B\mu_2)$ in
$\widehat M_{0,[6]}$.
\end{claim}
\begin{proof}
The displayed map gives a finite \'etale double cover of
$Y\setminus\{b\}$, followed by its one-puncture filling
into $Y$. On compactifications the two branch values are
$b$ and $\infty$. Every double cover with this labelled
target and these branch values is isomorphic to the
displayed one, uniquely up to its deck involution.
The family exists globally over $Y$, so the corresponding
$\mu_2$-gerbe is trivial: its stack is $Y\times B\mu_2$.
The source family gives $\phi$, with the $C_2$ factor
supplied by the global deck involution.

The source complexity is four, the target complexity is
one, and the degree is two, so the defect is two. The
single omitted ramification point over the interior
critical value $b$ already contributes all of that
defect. There are no further omitted interior points.
Thus any factorization of this composite into coverings
and fillings must first omit $b$ on the target so that
the degree-two factor is \'etale. Filling the source
ramification point before that factor would make it
ramified on the open. Apart from homeomorphisms and
trivial factors, the cover-then-fill pattern is forced.
Its diagram loops, with fixed labelled target, therefore
give exactly the right stabilizer $\operatorname{im}\phi$.

By Noohi's comparison, the \'etale group is the profinite
completion of $\pi_1(Y)\times C_2$. Its continuous image
in $\widehat M_{0,[6]}$ is compact and hence closed, and
the original image is dense in it. It is consequently
exactly $\overline{\operatorname{im}\phi}$.
No injectivity of $\phi$ or separability of its image
is required. This concrete braid calculation is
consistent with the general stabilizer proof of E.8.
\end{proof}

The Hausdorff clause in E.8 is deliberate. We do not
assert that the naive abstract quotient of finite words
in profinite paths has Hausdorff Hom spaces before this
separation. E3 is answered by the stated Hausdorff
presentation, with the precise closed-coset completion
and its algebraic comparison.


% Written from audit/F_E1_tame_construction.md and independently accepted by root.
\subsection{A repaired category using boundary-tame moduli paths}

The two changes in the next construction are explicit. In positive
characteristic we use paths tame at the boundary of the stable moduli
stack. In characteristic zero we restrict the covering arrows to those
whose Galois closures have order prime to the chosen residue
characteristic. Full characteristic-zero moduli markings remain allowed.
The obstructions E.1 and E.2 still apply to the original unrestricted
formulation.

\begin{claim}{E.10}{Proved}
Let $k$ be algebraically closed of characteristic $p>0$, and let
$\mathcal M=\mathcal M_{g,[n],k}$ be the stack of hyperbolic curves
with unordered punctures. Give its proper stable compactification
$\overline{\mathcal M}$ the divisorial logarithmic structure of its
boundary $D$. For $(m,p)=1$, put
\[
 \mathcal R_m=\sqrt[m]{(\overline{\mathcal M},D)},\qquad
 \mathcal P^{\mathrm{bt}}(\mathcal M)
   =\varprojlim_{(m,p)=1}\Pi_1^{\mathrm{et}}(\mathcal R_m),
\]
using geometric base points in $\mathcal M$. This is the full
boundary-tame fundamental groupoid. It is invariant under algebraically
closed field extension. For a smooth marked model over a strictly
henselian complete DVR of mixed characteristic $(0,p)$, it admits a
surjective specialization from the full geometric generic moduli
fundamental group, and the corresponding maps of path torsors are
surjective. Finally, the canonical map
\[
 \Aut_k(C)\longrightarrow
 \mathcal P^{\mathrm{bt}}(\mathcal M)(C,C)
\]
is injective for every hyperbolic curve $C/k$ of that type.
\end{claim}
\begin{proof}
We first justify the stack comparison being used. If $\mathcal X$ is
a proper finite-type Deligne--Mumford stack over a field or a complete
noetherian local ring, its diagonal is finite. Theorem~2.7 of
\citet{F_EHKV2001} gives a finite surjective representable map
$Y\to\mathcal X$ from a scheme. The scheme $Y$ and the next two
terms of its \v Cech nerve are proper schemes; the latter are schemes
because the diagonal is finite. Effective descent for finite
\emph{\'etale} covers along this map follows from SGA~1 IX,
Theorem~4.12, after an \emph{\'etale} atlas. It identifies a cover
of $\mathcal X$ with a finite \emph{\'etale} cover of $Y$, a
descent isomorphism on $Y\times_{\mathcal X}Y$, and its cocycle
on the triple product. Apply SGA~1 X, Corollary~1.8 to all three
proper schemes for algebraically closed field extension, or X,
Theorem~2.1 for passage from a complete local base to its closed
fibre. Full faithfulness carries the isomorphism and the cocycle
uniquely. Thus both finite-\emph{\'etale} comparison equivalences
hold for $\mathcal X$ \citep{E_SGA1}. This argument makes no
prime-to-$p$ assumption on the original stack inertia.

For completeness, the associated geometric specialization is
surjective when $\mathcal X$ is smooth proper over a complete DVR
$R$ with algebraically closed residue field. A connected finite
\emph{\'etale} special cover extends by the equivalence just
proved. The total cover is connected and regular. A connected
regular noetherian DM stack is irreducible, as may be checked on
\emph{\'etale} charts, so its generic fibre is connected.
Repeating after every finite extension of $R$ proves geometric
connectedness of that generic fibre. This is the Galois-category
criterion for surjectivity.

Knudsen's Theorem~2.7 supplies proper smoothness and relative normal
crossings for all stable pointed types; take the quotient by the
constant permutation group for unordered punctures
\citep[p.~179]{BC_Knudsen}. Locally, where the boundary parameters
are $x_1,\ldots,x_r$, the finite root stack is
\[
 \left[\Spec
 A[t_1,\ldots,t_r]/(t_1^m-x_1,\ldots,t_r^m-x_r)
       \big/\mu_m^r\right].
\]
The logarithmic construction glues these charts even when the
boundary has self-intersections or its local branches are permuted
\citep[Propositions~3.4 and 3.10]{F_TalpoVistoli2018}.
The stack is proper and smooth DM when $m$ is invertible. Its
finite representable \emph{\'etale} covers restrict fully faithfully
to the open moduli stack. Conversely, SGA~1 XIII, Proposition~5.2
extends a boundary-tame cover uniquely after extracting suitable
prime-to-$p$ roots of the boundary parameters. A common multiple
works on finitely many charts, and uniqueness of normalization
gives descent. This proves the groupoid description above. The
Kummer-flat theorem about infinite root stacks is not needed.

Apply the proper-stack comparisons to every $\mathcal R_m$.
They prove field-extension invariance. Over $R$, they give
surjections at each root level. The inverse-limit map is still
surjective: its compact image projects onto every level of a
directed inverse system and therefore is dense and closed.
Restriction of covers from a normal root stack to its dense open
also induces a surjection on fundamental groups. Combining these
maps proves the assertion for the full generic moduli group.
The assertion for path torsors follows by choosing a path and
using the group assertion. All ambiguities are the usual choices
of base points and paths.

It remains to prove the last assertion, which will exclude a
gerbe obstruction in E.11. Order the punctures; the ordering cover
extends as a finite \emph{\'etale} cover of the stable quotient
stack, even when $p\mid n!$. In genus zero an automorphism fixing
all these marks is identity. In positive genus add a full Jacobi
level at an odd prime $\ell\ne p$, $\ell\geq3$; in genus one
use an ordered mark as origin. These data kill automorphisms.
For genus at least two, use the faithful action on the Jacobian
in Deligne--Mumford Theorem~1.13 and Lemma~1.14. In genus one
use the faithful Tate-module action of origin-preserving elliptic
automorphisms. A finite-order matrix in
$1+\ell M_{2g}(\Z_\ell)$ is identity: the first nonzero
$\ell$-adic term of its powers, using $\ell$ odd, rules out both
prime-to-$\ell$ and $\ell$-power torsion. Thus a finite-order
automorphism fixing the full level acts trivially on the Tate
module and then on the curve
\citep[Theorem~1.13, Lemma~1.14, and \S5.14]{BC_DM}.

These level covers are boundary tame. On a regular chart of
the stable compactification, the universal curve is already
split semistable after strict localization. The prime-to-$p$
form of the monodromy computation in B.5 factors through tame
Kummer inertia, as in the proof of C.2. On homology modulo
$\ell$, a boundary generator acts by $1+N_e$, with
$N_e(x)=\langle x,\delta_e\rangle\delta_e$ and $N_e^2=0$.
Consequently $(1+N_e)^\ell=1$. At an intersection the
vanishing cycles are disjoint, so $N_eN_f=0$ and these
transvections generate an elementary abelian $\ell$-group.
Wild $p$-inertia is trivial. The level cover is therefore
tame at every boundary branch; the same argument applies
in the marked genus-one case. We use the prime-to-residue-
characteristic monodromy theorem underlying B.5 here, not a
claim about the full characteristic-$p$ geometric group
\citep[Theorem~4.8, Lemma~5.4]{BC_HMTopics}.
The action of $\Aut(C)$ on the fibre of the combined ordering
and level cover is faithful. Since this finite cover occurs in
the defining Galois category, its action detects every nonidentity
curve automorphism in $\mathcal P^{\mathrm{bt}}$.
\end{proof}

\begin{claim}{E.11}{Proved}
There is a category $\mathcal T_k^{\mathrm{bt}}$ with all
hyperbolic $k$-curves as objects, the groupoids
$\mathcal P^{\mathrm{bt}}$ as its marking groupoids, and all
actual finite \emph{\'etale} maps with prime-to-$p$ Galois
closure included faithfully. Every arrow factors as a marking
followed by such a cover. Covers transport along markings,
compatibly with composition. Algebraically closed field
extension induces an equivalence.

Write $c(X)=2g_X-2+n_X$. If $d=c(X)/c(Y)$ is not a positive
integer, then $\Hom(X,Y)$ is empty. Otherwise, for representatives
$u:V_u\to Y$ of the finitely many isomorphism classes of covers
of degree $d$, source type equal to that of $X$, and prime-to-$p$
Galois closure, there is a natural description
\[
 \Hom_{\mathcal T_k^{\mathrm{bt}}}(X,Y)
  =\coprod_u D_u\backslash
  \mathcal P^{\mathrm{bt}}(\mathcal M_{\mathrm{type}(X)})(X,V_u),
 \qquad D_u=\Aut_Y(V_u).
\]
These Hom spaces are compact Hausdorff, composition is continuous,
and $\Aut_{\mathcal T_k^{\mathrm{bt}}}(X)
=\pi_1^{\mathrm{bt}}(\mathcal M_{\mathrm{type}(X)},X)$.
\end{claim}
\begin{proof}
We give the path-lifting argument, including its inertia issue.
Let $u:V\to Y$ have Galois-closure group $G$ of order prime to
$p$. Use the stable Hurwitz stack with its Galois closure retained,
marking every puncture of $Y$ and every point above it. This is
the stabilized compactification of Bertin--Romagny, rather than
the generally nonnormal space of naive admissible covers.
Proposition~5.3, Definition~6.8, Theorem~6.9, Proposition~6.14,
Theorem~6.20 and Definitions~6.21--6.22 of \citet{E_BR2007}
give its properness, its source and target maps, and the
intermediate quotient constructions. All ramification points
are marked, so the nodal stabilizers are cyclic and do not
exchange branches. Their Theorem~5.5 gives the local
discriminant parameters
\[
                      \tau_i=t_i^{e_i},\qquad e_i\mid |G|,
\]
with the remaining deformation parameters unchanged.

Pull back a target root stack with $e_i\mid m$, and normalize.
The local branches are
$t_i=\zeta_i r_i^{m/e_i}$, where $\tau_i=r_i^m$.
Each branch has the same deformation ring as the root-stack
base. A relative automorphism fixing a branch fixes that full
deformation ring and extends uniquely over its universal
deformation. Its relative automorphism group is therefore the
deck group of the generic smooth cover. The extra boundary
automorphisms which permute the branches do not remain in the
relative stabilizer of a chosen normalization branch.
The identical computation for a non-Galois cover is
Bertin--Romagny Lemma~6.23 and Proposition~6.24. Consequently the
normalized morphism is, \emph{\'etale} locally on the target,
a finite disjoint union of finite \emph{\'etale} gerbes.
Equivalently it is such a gerbe over a representable finite
\emph{\'etale} cover of the target root stack. No
representability assertion is made for the Hurwitz discriminant
map itself; the nonrepresentability warning in Proposition~6.14
is essential here.

Because every $e_i$ is prime to $p$, the target roots so obtained
are cofinal among the prime-to-$p$ roots of the Hurwitz boundary.
The source map also acts on these tame groupoids: it extends
to the stable compactifications, and every source node parameter
pulls back to a monomial in Hurwitz node parameters times a unit.
The pullback of a prime-to-$p$ Kummer cover is again such a
cover. All inverse images of target marks have been retained;
this is what makes the intermediate quotient curves stable
and avoids the unmarked quotient problem noted in the source.

A finite \emph{\'etale} gerbe with fibre group $D$ has the
fundamental-group sequence
\[
       D\longrightarrow\pi_1(\mathcal H,u)
         \longrightarrow\pi_1(\mathcal Q,q)\longrightarrow1.
\]
The image of $D$ is normal and finite; this is Noohi
Proposition~10.4 and Corollary~10.5
\citep{F_Noohi2004}. The same assertion over a stack base is
proved by \emph{\'etale} descent: neutralize locally, check
it for $BD$, and observe that finite covers with trivial
relative-inertia action descend uniquely to $\mathcal Q$.
In general the left map need not be injective. Pass to the
cofinal root system first. The limit sequence remains exact:
surjectivity follows from the finite-intersection property
for compact lift sets; an element in the limit kernel comes
from a single element of the finite group $D$, again by the
finite-intersection property. The source map on this limit
sends $D$ to the canonical subgroup
$D\subset\Aut(V)\subset\pi_1^{\mathrm{bt}}(\mathcal M,V)$,
which is injective by E.10. Hence the left map is injective in
this application. This supplies the missing check that a
nontrivial gerbe transgression cannot kill deck transformations.

Thus a target tame path and a starting cover determine a unique
endpoint isomorphism class of covers and admit a lift in the
Hurwitz tame groupoid. Two lifts to the same endpoint differ by
a deck element. Their source paths therefore agree modulo the
action of that deck group.

Define an arrow by a pair $(u,\alpha)$, where $u:V\to Y$ is
one of the permitted covers and
$\alpha:X\rightsquigarrow V$ is a boundary-tame marking.
Identify $(u,\alpha)$ with $(u',h\alpha)$ when
$h:V\xrightarrow{\sim}V'$ is an isomorphism over $Y$.
For $(u,\alpha):X\to Y$ and $(v,\beta):Y\to Z$, write
$v:W\to Z$, transport $u$ along $\beta$ to
$u':V'\to W$, and let $\gamma:V\rightsquigarrow V'$ be its
source marking. Set
\[
                  (v,\beta)(u,\alpha)
                     =(vu',\gamma\alpha).
\]
The deck-kernel assertion proves independence of the lift and
of the chosen representatives. For associativity use the
Hurwitz stack of a chain of covers. A composite has Galois
closure contained in the corresponding wreath product, whose
order is still prime to $p$. Its common Galois closure and
the intermediate subgroups transport the chain simultaneously.
Their marked quotient models are stable by the result already
checked. Both bracketings give the same composite-cover map
and source path on that stack. Identities are $(\mathrm{id},
\mathrm{id})$.

Actual covers map to $(u,\mathrm{id})$. If two such maps became
equal, an isomorphism of their sources would have trivial
canonical marking. E.10 makes that isomorphism identity, so the
inclusion is faithful. Factorization and cover transport have
been built into the construction.

Tame Riemann--Hurwitz gives $c(V)=d\,c(Y)$. The maximal
prime-to-$p$ curve group is topologically finitely generated,
so it has only finitely many subgroups of a given index.
This proves the stated finite indexing of the Hom formula.
A path torsor is profinite and $D_u$ is finite; its quotient
is compact Hausdorff. On each component, endpoint transport
depends on a finite quotient, and after the endpoint is fixed
composition is induced by continuous maps of Hurwitz path
torsors and group multiplication. Composition is therefore
continuous. A degree-one cover is an isomorphism, giving the
asserted automorphism group.

Under $k\subset K$, E.3 identifies fixed-target covers, their
source types, and their deck groups. E.10 identifies the marking
torsors. The Hom formula proves full faithfulness. The moduli
stack of each hyperbolic type is geometrically connected, and
any $K$-curve has a type represented over $k$. A marking
connects it to such a base-changed curve. This proves essential
surjectivity, without asserting that the individual curve
descends as a scheme.
\end{proof}

\begin{claim}{E.12}{Proved}
Fix $p$, and let $\mathcal T^{\mathrm{cov}(p')}$ be the
characteristic-zero covering category with its full profinite
moduli markings, retaining precisely the covering arrows whose
Galois closures have order prime to $p$. After choosing good
reduction base models for the types and paths between type
base points, there is a specialization functor
\[
                 \mathcal T^{\mathrm{cov}(p')}
                     \longrightarrow \mathcal T_k^{\mathrm{bt}}.
\]
It preserves type and is surjective on each Hom space between
the chosen type base objects. On automorphism groups it is
\[
 \wh{\operatorname{Mod}}_{g,[n]}
       \twoheadrightarrow\pi_1^{\mathrm{bt}}(\mathcal M_{g,[n],k}).
\]
On Hom spaces it preserves the finite cover-component index
set and is the disjoint union of the surjections
\[
 D_u\backslash\operatorname{Path}_{\mathrm{et}}(X,V_u)
    \longrightarrow
 D_{u_s}\backslash\operatorname{Path}_{\mathrm{bt}}(X_s,V_{u_s}).
\]
The deck groups correspond under specialization. With compatible
path identifications and a group map $q:\Gamma\twoheadrightarrow
\Gamma_s$, the fibre above the coset of $q(\gamma)$ consists of
the left-$D_u$ classes in $D_u\ker(q)\gamma$.
\end{claim}
\begin{proof}
Take $R=W(k)$. For every type, smoothness of the pointed moduli
stack lifts a chosen special curve to a smooth marked model
over $R$; use a smooth chart through its point and completeness,
or formal existence for the proper curve. Fix one such model
per type. The full subcategories on their geometric generic
fibres and their special fibres are equivalent to the respective
full categories, because markings connect all objects of one
type. Characteristic-zero base-extension invariance identifies
the generic category with the stated one. One may use a common
algebraically closed overfield; no embedding of an arbitrarily
large field into $\C$ is required.

Let $u:V\to Y$ be a permitted cover of a chosen generic target.
By E.3 it corresponds to a special cover with the same degree,
source type and deck group. It has a smooth marked model over
the chosen target model after a finite extension of $R$:
lift the corresponding special cover, identify its geometric
generic fibre with $u$, and descend that identification to a
finite extension. Uniqueness of the marked stable model makes
this model and its deck action compatible with further
extensions.

Represent an arrow by $(u,\alpha)$, with $\alpha$ now a full
characteristic-zero moduli marking. Both endpoints of $\alpha$
have models after the preceding finite extension. E.10 specializes
the marking through the proper root stacks. Define
\[
                      (u,\alpha)\longmapsto
                         (u_s,\operatorname{sp}(\alpha)).
\]
The proper finite-\emph{\'etale} comparisons are compatible
with extension of the DVR. The separated stable moduli stack
gives the unique extensions of generic model isomorphisms.
These two facts give independence from extensions, representatives
and model identifications. The normalized root Hurwitz source
and target maps commute with specialization. Applying that
compatibility to the common-closure chain used in E.11 proves
compatibility with composition; hence this is a functor.

Degree and the marked source type are preserved by E.3. For
automorphisms, E.10 gives the surjection from the full generic
moduli group; the characteristic-zero comparison identifies it
with the full profinite mapping class group. More explicitly
the map first passes to the inverse limit of generic root
groups with indices prime to $p$, and then specializes. It
retains the target's finite interior $p$-monodromy; it is not
the maximal prime-to-$p$ quotient of the moduli group.

Finally apply the normal forms of E.11 and the fixed-target
cover equivalence E.3. Every cover component corresponds to
exactly one special component, with its same deck group.
Specialization of path torsors is surjective by E.10. It
therefore induces the displayed surjection on each component.
After identifying the torsors with their groups, two classes
have equal images exactly when their representatives differ
by a deck element and an element of $\ker(q)$, giving the
stated fibres. No injectivity of $q$ or faithfulness of this
specialization functor is asserted. E.2 proves that its
restriction on covering arrows cannot simply be removed.
\end{proof}


The former obligation E.U1 is superseded by E.10--E.12.
They prove the repaired construction with boundary-tame paths
retaining interior $p$-monodromy and the restricted-source
specialization. The literal full-path and unrestricted-source
requirements remain subject to E.1--E.2.

\subsection{The exact remaining arithmetic-image obligation}

\begin{claim}{E.U2}{Unresolved}
We do not identify the image of $G_\Q$ as compatible outer
automorphisms of all $\widehat M_{g,n}$ with all covering and
puncture-forgetting operations. E.4 supplies the published
genus-zero arithmetic subgroup characterization. Extending that
characterization to the specified entire category, with all
compatibilities and with the conjugacy ambiguity accounted for,
is the unproved step here. In particular we do not replace it by
the stronger and different equality $\widehat{\mathrm{GT}}=G_\Q$.
\end{claim}

E.8 resolves the earlier comparison obligation for the explicitly
Hausdorff algebraic presentation. It does not assert that the
unseparated abstract quotient of finite words in profinite paths
already has Hausdorff Hom spaces. That additional assertion is
unneeded for the completion theorem requested in E3.
